% Options for packages loaded elsewhere
\PassOptionsToPackage{unicode}{hyperref}
\PassOptionsToPackage{hyphens}{url}
\PassOptionsToPackage{dvipsnames,svgnames,x11names}{xcolor}
%
\documentclass[
  11pt,
]{article}
\usepackage{amsmath,amssymb}
\usepackage{iftex}
\ifPDFTeX
  \usepackage[T1]{fontenc}
  \usepackage[utf8]{inputenc}
  \usepackage{textcomp} % provide euro and other symbols
\else % if luatex or xetex
  \usepackage{unicode-math} % this also loads fontspec
  \defaultfontfeatures{Scale=MatchLowercase}
  \defaultfontfeatures[\rmfamily]{Ligatures=TeX,Scale=1}
\fi
\usepackage{lmodern}
\ifPDFTeX\else
  % xetex/luatex font selection
\fi
% Use upquote if available, for straight quotes in verbatim environments
\IfFileExists{upquote.sty}{\usepackage{upquote}}{}
\IfFileExists{microtype.sty}{% use microtype if available
  \usepackage[]{microtype}
  \UseMicrotypeSet[protrusion]{basicmath} % disable protrusion for tt fonts
}{}
\makeatletter
\@ifundefined{KOMAClassName}{% if non-KOMA class
  \IfFileExists{parskip.sty}{%
    \usepackage{parskip}
  }{% else
    \setlength{\parindent}{0pt}
    \setlength{\parskip}{6pt plus 2pt minus 1pt}}
}{% if KOMA class
  \KOMAoptions{parskip=half}}
\makeatother
\usepackage{xcolor}
\usepackage[margin=2.2cm]{geometry}
\setlength{\emergencystretch}{3em} % prevent overfull lines
\providecommand{\tightlist}{%
  \setlength{\itemsep}{0pt}\setlength{\parskip}{0pt}}
\setcounter{secnumdepth}{5}
\usepackage{lmodern}
\usepackage[T1]{fontenc}
\usepackage{amsmath,amssymb}
\usepackage{newunicodechar}
\newunicodechar{}{\ensuremath{\bar\nu}}
\newunicodechar{}{\ensuremath{\hat\nu}}
\newunicodechar{}{\ensuremath{\tilde\sigma}}
\newunicodechar{}{\ensuremath{\bar v}}
\newunicodechar{}{\ensuremath{\hat\sigma}}
\newunicodechar{}{\ensuremath{\bar X}}
\newunicodechar{}{\ensuremath{\hat X}}
\newunicodechar{}{\ensuremath{\tilde W}}
\newunicodechar{}{\ensuremath{\dot\alpha}}
\newunicodechar{}{\ensuremath{\dot\beta}}
\newunicodechar{}{\ensuremath{\hat f}}
\newunicodechar{}{\ensuremath{\bar t}}
\newunicodechar{}{\ensuremath{\hat V}}
\newunicodechar{Ĉ}{\ensuremath{\hat C}}
\newunicodechar{ĉ}{\ensuremath{\hat c}}
\newunicodechar{ġ}{\ensuremath{\dot g}}
\newunicodechar{ᵀ}{\ensuremath{{}^{\mathsf{T}}}}
\newunicodechar{ℝ}{\ensuremath{\mathbb{R}}}
\newunicodechar{−}{\ensuremath{-}}
\newunicodechar{²}{\ensuremath{{}^{2}}}
\newunicodechar{³}{\ensuremath{{}^{3}}}
\newunicodechar{¹}{\ensuremath{{}^{1}}}
\newunicodechar{⁴}{\ensuremath{{}^{4}}}
\newunicodechar{⁸}{\ensuremath{{}^{8}}}
\newunicodechar{⁺}{\ensuremath{{}^{+}}}
\newunicodechar{⁻}{\ensuremath{{}^{-}}}
\newunicodechar{₂}{\ensuremath{{}_{2}}}
\newunicodechar{½}{\ensuremath{\tfrac12}}
\newunicodechar{¼}{\ensuremath{\tfrac14}}
\newunicodechar{·}{\ensuremath{\cdot}}
\newunicodechar{×}{\ensuremath{\times}}
\newunicodechar{÷}{\ensuremath{\div}}
\newunicodechar{±}{\ensuremath{\pm}}
\newunicodechar{σ}{\ensuremath{\sigma}}
\newunicodechar{φ}{\ensuremath{\varphi}}
\newunicodechar{ν}{\ensuremath{\nu}}
\newunicodechar{τ}{\ensuremath{\tau}}
\newunicodechar{π}{\ensuremath{\pi}}
\newunicodechar{ρ}{\ensuremath{\rho}}
\newunicodechar{λ}{\ensuremath{\lambda}}
\newunicodechar{μ}{\ensuremath{\mu}}
\newunicodechar{κ}{\ensuremath{\kappa}}
\newunicodechar{ω}{\ensuremath{\omega}}
\newunicodechar{β}{\ensuremath{\beta}}
\newunicodechar{α}{\ensuremath{\alpha}}
\newunicodechar{γ}{\ensuremath{\gamma}}
\newunicodechar{θ}{\ensuremath{\theta}}
\newunicodechar{η}{\ensuremath{\eta}}
\newunicodechar{χ}{\ensuremath{\chi}}
\newunicodechar{ψ}{\ensuremath{\psi}}
\newunicodechar{ξ}{\ensuremath{\xi}}
\newunicodechar{δ}{\ensuremath{\delta}}
\newunicodechar{ε}{\ensuremath{\varepsilon}}
\newunicodechar{Σ}{\ensuremath{\Sigma}}
\newunicodechar{Δ}{\ensuremath{\Delta}}
\newunicodechar{Λ}{\ensuremath{\Lambda}}
\newunicodechar{Π}{\ensuremath{\Pi}}
\newunicodechar{Ξ}{\ensuremath{\Xi}}
\newunicodechar{Γ}{\ensuremath{\Gamma}}
\newunicodechar{Φ}{\ensuremath{\Phi}}
\newunicodechar{Θ}{\ensuremath{\Theta}}
\newunicodechar{Ψ}{\ensuremath{\Psi}}
\newunicodechar{∫}{\ensuremath{\int}}
\newunicodechar{∂}{\ensuremath{\partial}}
\newunicodechar{√}{\ensuremath{\surd}}
\newunicodechar{∞}{\ensuremath{\infty}}
\newunicodechar{≈}{\ensuremath{\approx}}
\newunicodechar{≥}{\ensuremath{\ge}}
\newunicodechar{≤}{\ensuremath{\le}}
\newunicodechar{≠}{\ensuremath{\ne}}
\newunicodechar{≡}{\ensuremath{\equiv}}
\newunicodechar{→}{\ensuremath{\to}}
\newunicodechar{←}{\ensuremath{\leftarrow}}
\newunicodechar{⇒}{\ensuremath{\Rightarrow}}
\newunicodechar{⇔}{\ensuremath{\Leftrightarrow}}
\newunicodechar{↦}{\ensuremath{\mapsto}}
\newunicodechar{∈}{\ensuremath{\in}}
\newunicodechar{∝}{\ensuremath{\propto}}
\newunicodechar{∏}{\ensuremath{\prod}}
\newunicodechar{⊥}{\ensuremath{\perp}}
\newunicodechar{‖}{\ensuremath{\|}}
\newunicodechar{⟨}{\ensuremath{\langle}}
\newunicodechar{⟩}{\ensuremath{\rangle}}
\newunicodechar{∩}{\ensuremath{\cap}}
\newunicodechar{⊆}{\ensuremath{\subseteq}}
\newunicodechar{…}{\ensuremath{\ldots}}
\newunicodechar{̄}{}
\newunicodechar{̂}{}
\newunicodechar{̃}{}
\newunicodechar{̇}{}
\ifLuaTeX
  \usepackage{selnolig}  % disable illegal ligatures
\fi
\IfFileExists{bookmark.sty}{\usepackage{bookmark}}{\usepackage{hyperref}}
\IfFileExists{xurl.sty}{\usepackage{xurl}}{} % add URL line breaks if available
\urlstyle{same}
\hypersetup{
  pdftitle={Exercise Set 2 --- Defense Sheet},
  pdfauthor={Computational Finance (WI4151), TU Delft --- Prof.~Fang Fang},
  colorlinks=true,
  linkcolor={blue},
  filecolor={Maroon},
  citecolor={Blue},
  urlcolor={Blue},
  pdfcreator={LaTeX via pandoc}}

\title{Exercise Set 2 --- Defense Sheet}
\author{Computational Finance (WI4151), TU Delft --- Prof.~Fang Fang}
\date{}

\begin{document}
\maketitle

{
\hypersetup{linkcolor=blue}
\setcounter{tocdepth}{3}
\tableofcontents
}
This set has two problems, each in its own file: - \textbf{Exercise 1
--- Least-Squares Monte Carlo (LSM / Longstaff--Schwartz):}
\texttt{exercise1\_def\ (2).py} - \textbf{Exercise 2 --- Finite
Difference (FTCS) for European and American puts:}
\texttt{exercise2\_def\ (2).py}

Line numbers below refer to those two files. Where a ``why'' requires my
own reasoning, it is left as a \(TODO (Bruno)\) for me to answer out
loud.

\begin{center}\rule{0.5\linewidth}{0.5pt}\end{center}

\hypertarget{the-task-in-my-own-words}{%
\section{The task (in my own words)}\label{the-task-in-my-own-words}}

\textbf{Exercise 1 --- LSM (Longstaff--Schwartz, 2001).} - Part (i):
reproduce the textbook 8-path numerical example from the paper. Strike
\(K = 1.10\), \(r = 6%
\), three yearly exercise dates (\(dt = 1\), \(T = 3\)). Print the
cash-flow matrix at each time, the in-the-money regression tables, the
fitted \(E[Y|X] = a0 + a1 X + a2 X^2\), the optimal early-exercise
decisions, the stopping rule, the option cash-flow matrix, and finally
the American and European put prices to 5 decimals. - Part (ii): replace
the 8 hard-coded paths with \(N = 100,000\) simulated GBM paths (50k +
50k antithetic), \(K = 40\), \(M = 50\) exercise dates per year, and
replicate the ``red box'' columns of Table 1 of the paper (Simulated
American, s.e., closed-form European, early-exercise value) for the 20
\((S0, sigma, T)\) combinations. Here the basis is the first three
weighted Laguerre polynomials evaluated at \texttt{x\ =\ S/K}.

\textbf{Exercise 2 --- Finite Difference (FTCS).} - Part (i): adapt a
provided \texttt{ftcs\_bs\_put.py} to price European puts via an
explicit (forward-time, central-space) scheme in log-price space
\(m = log S\), for the same 20 parameter combinations, and compare
against closed-form Black--Scholes. - Part (ii): turn the European FTCS
solver into an American one by adding the early-exercise (free-boundary)
projection \(U = max(U, intrinsic)\) after each time step. Compare
FD-American against the LSM-American of Exercise 1 as a correctness
check. - Part (iii): assemble the full ``outside the red box'' columns
of Table 1 (FD American, closed-form European, FD early-exercise value,
and the FD-minus-LSM difference in early-exercise value).

\textbf{TODO (Bruno):} In one sentence, why does the American put price
exceed the European put price here, and why is that \emph{not} true for
an American call on a non-dividend stock?

\begin{center}\rule{0.5\linewidth}{0.5pt}\end{center}

\hypertarget{my-approach-and-why}{%
\section{My approach and why}\label{my-approach-and-why}}

\textbf{Exercise 1 --- LSM.} - Backward induction over exercise dates.
At each date, regress discounted realized future cash flows \(Y\) on a
basis in the current spot \(X\), but \textbf{only over in-the-money
paths}, then compare immediate exercise against the fitted continuation
value. - Part (i) uses a degree-2 monomial basis \({1, X, X^2}\) via an
sklearn pipeline (\texttt{exercise1\_def\ (2).py:107}). - Part (ii) uses
the first three weighted Laguerre polynomials in \texttt{x\ =\ S/K}
(\texttt{exercise1\_def\ (2).py:185-191}). - Variance reduction by
antithetic paths (\texttt{exercise1\_def\ (2).py:177-178}).

\textbf{TODO (Bruno):} Why regress only on in-the-money paths rather
than all paths? (Out-of-the-money continuation values are irrelevant to
the exercise decision --- but say \emph{why} in your own words and tie
it to the paper.)

\textbf{TODO (Bruno):} Why monomials \({1, X, X^2}\) in Part (i) but
weighted Laguerre in Part (ii)? Was this just ``match the paper,'' or
numerical conditioning? Be ready to defend both.

\textbf{TODO (Bruno):} Antithetic variates: why does pairing \(Z\) with
\(-Z\) reduce variance for this payoff, and what property of the
estimator makes the variance reduction work? State the covariance
argument.

\textbf{Exercise 2 --- FTCS.} - Work in log-price \(m = log S\) so the
Black--Scholes PDE has constant coefficients, build the explicit
evolution matrix \(F\) once per parameter set, and march forward in
transformed time (time-to-maturity). - American: project onto the payoff
after each step (\texttt{exercise2\_def\ (2).py:105}).

\textbf{TODO (Bruno):} Why transform to log-price \(m = log S\) before
discretizing instead of working directly in \(S\)? (Constant
coefficients / uniform grid --- but explain the consequence for the
matrices \(T1\), \(T2\).)

\textbf{TODO (Bruno):} FTCS is \emph{explicit}. Why is that an
acceptable choice here, and what stability condition are you implicitly
relying on with \(Nt = 40000*T\)? (See ``hardest question'' below ---
this is the danger zone.)

\begin{center}\rule{0.5\linewidth}{0.5pt}\end{center}

\hypertarget{key-code-sections-file-what-it-does-why-i-wrote-it-that-way}{%
\section{Key code sections (file + what it does + why I wrote it that
way)}\label{key-code-sections-file-what-it-does-why-i-wrote-it-that-way}}

\hypertarget{exercise-1-exercise1_def-2.py}{%
\subsection{\texorpdfstring{Exercise 1 ---
\texttt{exercise1\_def\ (2).py}}{Exercise 1 --- exercise1\_def (2).py}}\label{exercise-1-exercise1_def-2.py}}

\begin{itemize}
\tightlist
\item
  \textbf{\texttt{:11-24}} --- Part (i) parameters and the 8 hard-coded
  paths from the paper. \(C\) (cash-flow matrix) initialized to zeros at
  \texttt{:27}, shape \((8, 3)\).
\item
  \textbf{\texttt{:32-86}} --- three printing helpers
  (\texttt{print\_formatted\_table}, \texttt{print\_regression\_table},
  \(optimal_table\)) that produce the per-time tables shown in the
  report. Pure presentation, no pricing logic.
\item
  \textbf{\texttt{:89} \(LSQM(S,C,dt,T,K,r)\)} --- the Part (i) LSM
  engine.

  \begin{itemize}
  \tightlist
  \item
    \texttt{:91} terminal payoff \(C[:,T-1] = max(K - S[:,T], 0)\).
  \item
    \texttt{:93} \(C_european\) keeps a copy of the terminal payoff for
    the European price later.
  \item
    \texttt{:96} backward loop \(i = T-2 ... 0\).
  \item
    \texttt{:98} regressor \(X = S[:,i+1]\) masked by \(<K\) (so OTM
    entries are zeroed before the ITM filter), reshaped to 2D.
  \item
    \texttt{:99-101} builds \(Y\) = sum of \emph{all} future cash flows
    discounted back to time \(i+1\) using
    \(discount_factors = exp(-r*dt*(u-i))\) with
    \(u = arange(i+1, ncols)\).
  \item
    \texttt{:102-104} ITM mask \(S[:,i+1] < K\); restrict \(X\), \(Y\)
    to ITM.
  \item
    \texttt{:107-108} degree-2 polynomial + linear regression pipeline,
    fit on ITM data.
  \item
    \texttt{:111-112} extracts intercept \(a0\) and coefficients
    \(a1, a2\).
  \item
    \texttt{:117-120} continuation = model prediction on ITM paths;
    exercise = \(K - S\) on ITM paths.
  \item
    \texttt{:122-127} exercise rule: if \(exercise > continuation\), set
    \(C[j,i]\) to exercise value and \textbf{zero out all later cash
    flows} \(C[j,i+1:] = 0\); else \(C[j,i] = 0\).
  \item
    \texttt{:131} stopping matrix from \(C != 0\).
  \item
    \texttt{:147-152} American price: discount each path's cash flows by
    \(exp(-r*dt*t)\) for \(t=1..T\), sum per path, average.
  \item
    \texttt{:156} European price: \(mean(C_european * exp(-r*dt*T))\).
  \end{itemize}
\item
  \textbf{\texttt{:167-173}} --- Part (ii) parameters (\(K=40\),
  \(M=50\), \(N=100000\)).
\item
  \textbf{\texttt{:175-183} \(simulate_paths\)} --- GBM exact-step
  simulation; \texttt{:177} draws \texttt{N//2} normals, \texttt{:178}
  stacks \(-Z\) for antithetics; \texttt{:181-182} exact log-Euler
  update
  \(S[:,i+1] = S[:,i]*exp((r-0.5 sigma^2)dt + sigma sqrt(dt) Z)\).
\item
  \textbf{\texttt{:185-191} \(laguerre_basis\)} --- \texttt{x=S/K};
  \texttt{L0=exp(-x/2)}, \texttt{L1=exp(-x/2)(1-x)},
  \texttt{L2=exp(-x/2)(1-2x+x\^{}2/2)}; stacked column-wise.
\item
  \textbf{\texttt{:193-196} \texttt{black\_scholes\_put}} ---
  closed-form European put.
\item
  \textbf{\texttt{:198-235} \(LSM_pricing\)} --- Part (ii) LSM engine;
  same backward structure as \(LSQM\) but basis is \(laguerre_basis\)
  (\texttt{:212}) and \(LinearRegression()\) with default intercept
  (\texttt{:215}). Returns per-path discounted cash flows.
\item
  \textbf{\texttt{:237-267}} --- loop over the 20 \((S0, sigma, T)\)
  combos: \(M_total = M*T\), \texttt{dt\ =\ 1/M}, simulate, price,
  compute \(american = mean(cashflow)\),
  \texttt{european\ =\ black\_scholes\_put(...)},
  \(ee = |american - european|\) (\texttt{:257}),
  \texttt{se\ =\ std(cashflow)/sqrt(N)} (\texttt{:258}).
\end{itemize}

\textbf{TODO (Bruno):} \texttt{:91} and \texttt{:203} index \(S[:, T]\)
/ \(S[:, M_total]\) for the terminal payoff while the cash-flow matrix
\(C\) has only \(T\) (resp. \(M_total\)) columns indexed \(0..T-1\).
Explain the off-by-one convention: \(S\) has one more column (the
\(t=0\) spot) than \(C\). Be ready to point at exactly which column of
\(S\) is ``today'' and which is maturity.

\textbf{TODO (Bruno):} \texttt{:257} uses
\texttt{ee\ =\ np.abs(american\ -\ european)} (absolute value) in the
Part (ii) block of \texttt{exercise1\_def}, but the copy in
\texttt{exercise2\_def\ (2).py:223} uses the \emph{signed}
\texttt{ee\ =\ american\_value\ -\ european\_value}. Which one matches
the paper's ``early exercise value'' column, and is the absolute value
ever wrong (can LSM-American dip below European by noise)?

\textbf{TODO (Bruno):} Standard error
\texttt{se\ =\ std(cashflow)/sqrt(N)} at \texttt{:258}. With antithetic
sampling the \(N\) draws are \textbf{not independent} (each path is
paired with its antithetic). Is dividing by \(sqrt(N)\) the correct
s.e., or does it understate/overstate it? Defend your formula.

\hypertarget{exercise-2-exercise2_def-2.py}{%
\subsection{\texorpdfstring{Exercise 2 ---
\texttt{exercise2\_def\ (2).py}}{Exercise 2 --- exercise2\_def (2).py}}\label{exercise-2-exercise2_def-2.py}}

\begin{itemize}
\tightlist
\item
  \textbf{\texttt{:5-14}} --- parameters \(K=40\), \(r=0.06\), and the
  20 \((S0, sigma, T)\) combos.
\item
  \textbf{\texttt{:18-64}} --- \textbf{Part (i) European FTCS} loop:

  \begin{itemize}
  \tightlist
  \item
    \texttt{:24-26} closed-form BS put for reference.
  \item
    \texttt{:29} domain length \(L = 2*log(K) + 2\); \texttt{:30}
    \(Nx = 1000\); \texttt{:31} \(Nt = 40000*T\); \texttt{:32-33} steps
    \texttt{h\ =\ L/Nx}, \texttt{k\ =\ T/Nt}.
  \item
    \texttt{:35-36} central-difference matrices: \(T1\) (first
    derivative, \texttt{+1/-1} off-diagonals), \(T2\) (second
    derivative, \(-2\) on diagonal, \(1\) off-diagonals), size
    \((Nx-1)\).
  \item
    \texttt{:38-40} explicit evolution matrix
    \texttt{F\ =\ (1\ -\ r\ k)\ I\ +\ 0.5\ k\ sigma\^{}2\ /\ h\^{}2\ *\ T2\ +\ k\ (r\ -\ 0.5\ sigma\^{}2)/(2h)\ *\ T1}.
  \item
    \texttt{:42-43} grid \(mvec\), centered and shifted by \(+log(K)\).
  \item
    \texttt{:45-46} initial condition \(U[:,0] = max(K - exp(mvec), 0)\)
    (payoff as the time-0 condition in time-to-maturity).
  \item
    \texttt{:48-52} forward march; \texttt{:51} boundary correction term
    \(p[0]\) injected at the low-\(m\) (deep ITM) boundary node, scaled
    by \(K*exp(-r*time2mat)\).
  \item
    \texttt{:54} \texttt{np.interp} to read off the price at
    \(log(S0)\).
  \end{itemize}
\item
  \textbf{\texttt{:66-118}} --- \textbf{Part (ii) American FTCS} loop:
  identical to Part (i) except \(intrinsic = max(K - exp(mvec), 0)\)
  (\texttt{:95}) and the single added projection line \texttt{:105}
  \texttt{U{[}:,i+1{]}\ =\ np.maximum(U{[}:,i+1{]},\ intrinsic)}.
  \texttt{ee\ =\ ftcs\_american\ -\ bs\_put} (\texttt{:108}).
\item
  \textbf{\texttt{:120-204}} --- re-imports and re-runs the
  \textbf{Exercise 1 Part (ii) LSM} code (paths, Laguerre basis, BS put,
  \(LSM_pricing\)) to regenerate \(results_lsmc\) for the cross-check.
\item
  \textbf{\texttt{:235-248}} --- correctness test table:
  \(|FD American - LSM American|\) per row.
\item
  \textbf{\texttt{:250-261}} --- full ``outside the red box'' table
  including the signed \(FD EE - LSM EE\) difference.
\end{itemize}

\textbf{TODO (Bruno):} \texttt{:29} \(L = 2*log(K) + 2\). Why this
width, and why center the grid at \(log(K)\) (\texttt{:43})? What goes
wrong at short maturity / high vol if the domain is too narrow --- and
what is the analogy with the COS truncation-range \([a,b]\) issue I hit
in Assignment 2?

\textbf{TODO (Bruno):} \texttt{:51} the boundary term \(p[0]\). Explain
in words what Dirichlet boundary condition this enforces at the
deep-in-the-money edge, why it carries the factor
\(K*exp(-r*time2mat)\), and what happens at the \emph{other}
(high-\(S\)) boundary where no correction is added.

\textbf{TODO (Bruno):} \texttt{:105} is described in the report as ``the
single line that changes.'' Explain why projecting
\(U = max(U, intrinsic)\) after each explicit step is a valid
(operator-splitting / explicit-penalty) treatment of the American free
boundary, and what its accuracy cost is versus a proper LCP/PSOR solve.

\begin{center}\rule{0.5\linewidth}{0.5pt}\end{center}

\hypertarget{design-decisions-i-must-justify}{%
\section{Design decisions I must
justify}\label{design-decisions-i-must-justify}}

\begin{enumerate}
\def\labelenumi{\arabic{enumi}.}
\tightlist
\item
  \textbf{ITM-only regression} (\texttt{:103-104}, \texttt{:211-213}).
  Factual: the mask is \(S[:,i+1] < K\). \textbf{TODO (Bruno):} justify
  why this is the Longstaff--Schwartz prescription and not a shortcut.
\item
  \textbf{Basis choice:} monomials degree 2 in Part (i) (\texttt{:107})
  vs weighted Laguerre in Part (ii) (\texttt{:185-191}). Factual: the
  report says normalization \texttt{x=S/K} avoids underflow in
  \texttt{exp(-x/2)} for \texttt{S\ in\ {[}36,44{]}}. \textbf{TODO
  (Bruno):} confirm you can derive the underflow argument and say why
  monomials were fine for the small \(K=1.10\) example.
\item
  \textbf{Antithetic variates} (\texttt{:177-178}). \textbf{TODO
  (Bruno):} justify and quantify expected variance reduction.
\item
  \textbf{Exact GBM step} (\texttt{:182}) rather than Euler.
  \textbf{TODO (Bruno):} why is the log update exact for GBM, and does
  discretization bias still enter the American price through the
  discrete exercise dates regardless?
\item
  \textbf{\(M = 50\) exercise dates/year as a proxy for continuous
  American exercise} (\texttt{:172}, \texttt{:251}). \textbf{TODO
  (Bruno):} why is a Bermudan with 50/100 dates a good stand-in for the
  true American, and which direction is the bias?
\item
  \textbf{Explicit FTCS with \(Nt = 40000*T\)}
  (\texttt{exercise2\_def\ (2).py:31}). \textbf{TODO (Bruno):} show that
  this choice keeps the scheme inside the explicit stability region
  (compute the CFL-type bound; see hardest question).
\item
  \textbf{\texttt{np.interp} to extract the price at \(log(S0)\)}
  (\texttt{:54}, \texttt{:107}). Factual: linear interpolation between
  grid nodes. \textbf{TODO (Bruno):} is linear interpolation accurate
  enough given \(h\), or should it be higher order?
\end{enumerate}

\begin{center}\rule{0.5\linewidth}{0.5pt}\end{center}

\hypertarget{results-and-what-they-mean}{%
\section{Results and what they mean}\label{results-and-what-they-mean}}

\textbf{Exercise 1, Part (i)} (from the report): American put price
\textbf{0.11443}, European put price \textbf{0.05638} on the same 8
paths. The fitted regressions match the paper: at \(t=2\),
\(E[Y|X] = -1.06999 + 2.98341 X - 1.81358 X^2\); at \(t=1\),
\(E[Y|X] = 2.03751 - 3.33544 X + 1.35646 X^2\). Stopping rule and option
cash-flow matrix reproduce the paper's table.

\textbf{Exercise 1, Part (ii)} (report Fig. 13): the 20-row replication
of the red-box columns; e.g.~\(S=36, sigma=0.2, T=1\) gives Simulated
American \textbf{4.4714 (s.e. 0.0094)}, closed-form European
\textbf{3.8443}, early-exercise value \textbf{0.6271}. s.e. all in the
0.006--0.023 range.

\textbf{Exercise 2, Part (i)} (report Fig. 14): FD-European vs
BS-European agree to \textasciitilde1e-4 across all 20 rows (largest
\(|Difference| = 0.0008\) at \(S=40, sigma=0.2, T=1\)).

\textbf{Exercise 2, Part (ii)} (report Fig. 16): FD-American prices, all
above the corresponding European, with early-exercise values matching
the LSM ones in sign and magnitude.

\textbf{Exercise 2 cross-check} (report Fig. 17):
\(|FD American - LSM American|\) is below 0.05 for all 20 rows (max
\textasciitilde0.016), and the Part (iii) \(FD EE - LSM EE\) differences
are small and alternate in sign.

\textbf{TODO (Bruno):} The FD-European prices are \emph{systematically
slightly below} BS (every \(Difference\) is negative, \textasciitilde{}
-1e-4 to -8e-4). Is that a bias of the explicit scheme, the truncated
domain, or the interpolation? Pick one and defend it --- Fang will ask
why the sign is consistent.

\textbf{TODO (Bruno):} The FD-American vs LSM-American differences are
\textasciitilde0.005--0.016, an order of magnitude larger than the
FD-European vs BS error (\textasciitilde1e-4). Why is the American
comparison so much looser? (Hint: LSM bias + 50-date Bermudan +
regression noise vs a near-exact PDE.)

\begin{center}\rule{0.5\linewidth}{0.5pt}\end{center}

\hypertarget{the-hardest-question-fang-could-ask-here-and-my-answer}{%
\section{The hardest question Fang could ask here --- and my
answer}\label{the-hardest-question-fang-could-ask-here-and-my-answer}}

\textbf{Likely hardest: ``Your FTCS scheme is explicit. Is it stable for
\(Nx = 1000\), \(Nt = 40000*T\)? Derive the condition and check it for
your worst-case \(sigma = 0.4\).''}

Factual setup from the code: in log-space the diffusion coefficient is
\(0.5*sigma^2\), \texttt{h\ =\ L/Nx} with
\(L = 2 log K + 2 = 2 log 40 + 2 ≈ 9.376\), so \(h ≈ 0.00938\);
\texttt{k\ =\ T/Nt\ =\ 1/40000\ =\ 2.5e-5} per unit time.

The explicit-scheme stability constraint for the diffusion part is
roughly \texttt{0.5*sigma\^{}2\ *\ k\ /\ h\^{}2\ \textless{}=\ 1/2},
i.e.~\texttt{k\ \textless{}=\ h\^{}2\ /\ sigma\^{}2}.

\textbf{TODO (Bruno):} Plug in the numbers:
\texttt{h\^{}2\ /\ sigma\^{}2\ =\ (0.00938)\^{}2\ /\ 0.16\ ≈\ 5.5e-4},
and \(k = 2.5e-5 << 5.5e-4\), so the diffusion number
\(≈ 0.045 << 0.5\). State this out loud, confirm the scheme is
comfortably stable, and explain \emph{why the author picked Nt so large}
(was it stability, or matching the paper's accuracy?). Also be ready to
say what the \((1 - r k)\) reaction term and the \(T1\) advection term
contribute to the stability bound --- and whether the advection
(first-derivative) term can cause spurious oscillations even when the
diffusion bound is met.

\textbf{Backup hardest (LSM): ``Is your LSM price biased high or low,
and why?''} \textbf{TODO (Bruno):} Standard result --- using the
\emph{same} paths to estimate the continuation regression and to value
the option introduces a look-ahead/in-sample bias in the continuation
estimate, so the exercise decision uses noisy information. State the
direction of the resulting price bias and how an out-of-sample /
two-pass scheme would fix it.

\begin{center}\rule{0.5\linewidth}{0.5pt}\end{center}

\hypertarget{show-me-in-your-code-where-you-do-x-anticipated-spots}{%
\section{``Show me in your code where you do X'' --- anticipated
spots}\label{show-me-in-your-code-where-you-do-x-anticipated-spots}}

\begin{itemize}
\tightlist
\item
  \textbf{``\ldots filter to in-the-money paths before regressing.''}
  \texttt{exercise1\_def\ (2).py:102-104}
  (\(itm_indexes = S[:,i+1] < K\)), and \texttt{:211-213} for Part (ii).
\item
  \textbf{``\ldots build the discounted continuation target Y.''}
  \texttt{:99-101} (Part i) and \texttt{:207-209} / \texttt{:182-184}
  (Part ii).
\item
  \textbf{``\ldots fit the continuation-value regression and read off
  the coefficients.''} \texttt{:107-112} (monomial pipeline) and
  \texttt{:215-217} (Laguerre).
\item
  \textbf{``\ldots make the exercise-vs-continue decision and zero out
  future cash flows.''} \texttt{:122-127} (Part i), \texttt{:224-229} /
  \texttt{:194-199} (Part ii).
\item
  \textbf{``\ldots compute the American and European prices.''}
  \texttt{:147-157} (Part i American + European).
\item
  \textbf{``\ldots generate antithetic paths.''} \texttt{:177-178}.
\item
  \textbf{``\ldots evaluate the weighted Laguerre basis.''}
  \texttt{:185-191}.
\item
  \textbf{``\ldots the closed-form Black--Scholes put.''}
  \texttt{:193-196} (Ex.1), \texttt{exercise2\_def\ (2).py:24-26} and
  \texttt{:74-76}.
\item
  \textbf{``\ldots assemble the explicit FTCS evolution matrix F.''}
  \texttt{exercise2\_def\ (2).py:35-40} (and \texttt{:85-90}).
\item
  \textbf{``\ldots set the initial/boundary conditions of the PDE.''}
  initial: \texttt{:45-46}; boundary term: \texttt{:51} (and
  \texttt{:103}).
\item
  \textbf{``\ldots the one line that makes it American.''}
  \texttt{exercise2\_def\ (2).py:105}
  (\texttt{U\ =\ np.maximum(U,\ intrinsic)}).
\item
  \textbf{``\ldots read the price off the grid at S0.''} \texttt{:54}
  and \texttt{:107}
  (\texttt{np.interp(np.log(S0),\ mvec,\ U{[}:,\ Nt{]})}).
\item
  \textbf{``\ldots cross-check FD against LSM.''} \texttt{:235-248}
  (\(|FD American - LSM American|\)).
\end{itemize}

\end{document}
